\section{Combat}\label{sec:8.0}

\ruleslabel{General Rule}

\begin{generalrule}

Each unit in the game possesses a Combat Strength (or simply ‘‘Strength’’) expressed in points, which is employed when attacking or defending. Only opposing units occupying corresponding boxes, triangles, or circles may participate in combat. Combat is mandatory between opposing units in corresponding Objective Display boxes and voluntary between opposing units in corresponding Mountain or Hide-away triangles or circles. Only the Phasing Player may initiate combats. If a combat is initiated, the Phasing Player is the attacker and the non-Phasing Player is the defender. Combat results will obligate the losing Player to lose Strength Points and/or retreat his units. The outcome of a combat may be influenced by terrain.

\end{generalrule}

\ruleslabel{Cases}

\subsection{How Combat Occurs}\label{sec:8.1}

\case{8.11} Combat may occur only under the following circumstances:

\textbf{A.} Opposing units occupy corresponding boxes of an Objective Display (see 5.3). In this situation, the Phasing Player must initiate combat during his Combat Phase;

\textbf{B.} Opposing units occupy corresponding Mountain triangles on a given Zone Display (see 5.2). Combat is voluntary in this situation; i.e., the Phasing Player is not obligated to initiate it during his Combat Phase;

\textbf{C.} Opposing units occupy corresponding Hide-away circles on a given Zone Display (see 5.2). Combat is voluntary in this situation; i.e., the Phasing Player is not obligated to initiate it during his Combat Phase. Note: Axis units may only enter Hide-away boxes on Anti-Guerrilla Operations; see 8.44.

\case{8.12} If the Phasing Player has initiated a combat within an Objective or Zone Display, all of both Players’ units occupying these corresponding boxes, triangles, or circles must participate. Exceptions: See 7.33, 8.25, and 8.48. No units may be voluntarily withheld from this combat.

\case{8.13} No unit may attack more than once per Combat Phase and no unit may be attacked more than once per Combat Phase. If an Axis stack in a Mountain or Hide-away has Partisan and Chetnik units in corresponding Mountain triangles or circles, this stack may attack only the Partisan or the Chetnik stack in a given Combat Phase, never both.

\subsection{How To Resolve Combat}\label{sec:8.2}

\case{8.21} If an attack has been declared by the Phasing Player, it is resolved in the following manner:

\textbf{A.} The Phasing and the non-Phasing Player total the Strength Points of their respective participating units;

\textbf{B.} The Phasing Player’s total is compared to the non-Phasing Player’s total, and this comparison is stated as a probability ratio: attacker’s Strength to defender’s Strength. This ratio is rounded down in favor of the defender to conform to the simplified ratio columns found on the Combat Results Table (CRT). Note: Terrain may affect the Strength of the defender (see 8.3);

\textbf{C.} The Phasing Player rolls a single die and crossreferences this die roll with the appropriate ratio column. The result is applied immediately (see 8.22).

\case{8.22} Combat results are read as follows:

\textbf{A.} A (attacker) or D (defender) followed by a number: The affected Player loses a unit(s) equal in Strength Points to the numerical result. If it is impossible for this Player to equal exactly the numerical result in Strength Points, he must eliminate as many Strength Points as possible from the affected stack without exceeding the numerical result. If it is impossible to lose any units that would yield a total of eliminated Strength Points less than the numerical result, nothing is lost. Regardless of how capable a Player’s force is of fulfilling a loss result, this force must retreat after combat (see 8.24);

\textbf{B.} DE or AE: The defender’s or attacker’s force is eliminated from play entirely.

Example: A Partisan brigade (Strength of 4) receives a D2 result. It is impossible for the Yugoslav Player to lose the brigade without suffering a Strength Point loss of less than 2. Thus, no loss is taken, but the brigade must retreat.

\case{8.23} Yugoslav brigades and divisions may never be broken down in order to fulfill combat loss results.

\case{8.24} Retreats are conducted as follows:

\textbf{A.} Yugoslav units always retreat into the Hide-away circle of the Zone they currently occupy. However, in all adverse results suffered in Axis Anti-Guerrilla Operations (see 8.4), Yugoslav units must retreat into the Hide-away of any adjacent Zone (Axis Player’s choice, subject to the restrictions of Case 6.6);

\textbf{B.} When engaged in a combat entirely against guerrilla units in a Mountain triangle, Axis units incurring a retreat result are not obligated to retreat.

Note: Axis units on Anti-Guerrilla Operations are never obligated to retreat.

\case{8.25} Units that are retreated may not participate in combat again for the duration of the Combat Phase, even if the box, triangle, or circle they are retreated into possesses units that are attacking or defending. If units retreat into a box, triangle, or circle that subsequently suffers an adverse combat result in the same Phase, these units are immediately eliminated. (Since stacking limitations apply only at the end of a Friendly Movement Phase, units may be retreated into a location in excess of these limitations.)

\case{8.26} Units may never ‘‘advance’’ after combat.

\case{8.27} The Phasing Player conducts combats in any order he wishes.

\case{8.28} If units are unable to retreat due to the restrictions of Case 6.6, they are eliminated.

\subsection{Effects Of Terrain On Combat}\label{sec:8.3}

\case{8.31} Yugoslav units occupying Hide-away circles are doubled in Combat Strength when defending.

\case{8.32} The numerical loss result following the letters A or D on the CRT is modified according to the following schedule:

\textbf{A.} It is tripled if the combat is taking place in a city (i.e., a D1 would become a D3);

\textbf{B.} It is doubled if the combat is taking place in a town or market town;

\textbf{C.} It is doubled in all combats taking place during an Anti-Guerrilla Operation (see 8.4).

\subsection{Anti-Guerrilla Operations}\label{sec:8.4}

Between Game-Turns 3 and 14 (inclusive), the Axis Player may conduct up to two
Anti-Guerrilla Operations per Game-Turn. However, the number of Anti-Guerrilla
Operations that may be performed during the course of the game is limited (see
15.3). Anti-Guerrilla Operations allow Axis units to enter the Hide-away circles
of the Zones in which the operations are being conducted and permit more
favorable attacks to be performed by the Axis Player.

\case{8.41} At the beginning of each Axis Anti-Guerrilla Operations Phase between Game-Turns 3 and 14 (inclusive), the Axis Player must state if he is performing none, one, or two operations in the current Game-Turn. If an operation(s) is declared, the Axis Player secretly writes down the name of a single Zone in which this operation will be conducted. (If two operations are declared, two Zones are listed — but the same Zone may not be chosen twice in the same Game-Turn). If no operations are declared, the rest of this Phase is skipped.

\case{8.42} For each declared operation, the Axis Player may choose up to 7 divisions (or their equivalents; see 6.22) during the Planning Segment to participate. These choices are made openly but may only include units that would normally be able to reach the Zone through normal movement in a hypothetical Movement Phase. In addition, Occupation Zone movement restrictions apply fully to Anti-Guerrilla Operations (see 6.6). After these units have been chosen, they are removed from the map and placed aside. (Note: See 10.23.)

\case{8.43} During the Yugoslav Reaction Segment, the Yugoslav Player rolls a single die (even if two operations have been declared). The result is the number of Yugoslav units that may be immediately permitted a ‘‘bonus’’ Movement Phase, subject to all normal movement rules. The units chosen to make this bonus move may be picked from any area on the map, although the Yugoslav

Player is not obligated to move any units if he does not wish to. The movement
of Tito counts against the Yugoslav Player’s allotment of bonus moves. Units
chosen to make a bonus move may move normally in the ensuing Yugoslav
Player-Turn.

\case{8.44} During the Deployment Segment, all Axis units participating in an Anti-Guerrilla Operation must be placed in the Mountain triangle and/or Hide-away circle of the Zone Display declared as the object of this operation. Participating units may be placed in these two locations in any combination the Axis Player desires — all may be placed in one or the other, or in any other conceivable combination. If two operations have been declared in the current Game-Turn, the Axis Player must resolve the first fully before proceeding to deploy the participating units of the second (i.e., performing Steps C and D of this Phase).

\case{8.45} During the Combat Segment, Axis units in a Mountain triangle or Hide-away circle may attack Yugoslav units in corresponding triangles and circles subject to normal combat rules. If no Yugoslav units are present in the corresponding locations, the Axis units are immediately removed from the map (see 8.47).

\case{8.46} The following special rules apply to combat during Anti-Guerrilla Operations:

\textbf{A.} The ratio is shifted two columns to the right (i.e., a 3-1 becomes a 5-1);

\textbf{B.} All numbered results on the CRT are doubled (i.e., a D1 becomes a D2);

\textbf{C.} If the Yugoslav Player is forced to retreat, his units are retreated into the Hide-away circle of an adjacent Zone;

\textbf{D.} Axis units are never obligated to retreat.

\case{8.47} After each combat is resolved, participating Axis units are
immediately removed from the map. During the Anti-Guerrilla Operations
Redeployment Phase (of the Terminal Stage), these units are placed back on the
map. At this time, the Axis Player may place each unit in any Axis box or
triangle of the Zone in which the operation was conducted---Exception: see
6.64B---including those with corresponding Enemy units.

\case{8.48} Axis units that are not participating in an operation may not be included in any attacks during this Phase, even if they occupy the same location as participating Axis units.

\case{8.49} Immediately before the start of the game, the Axis Player must determine the number of Anti-Guerrilla Operations that he is limited to for the duration of the game (see 15.3).

\subsection{Intrinsic Defense Strengths Of Axis Towns And Cities}\label{sec:8.5}

Trieste, Pola, Fiume, Zara, and Belgrade possess intrinsic Defense Strengths which are printed in parentheses within the Axis box of these Objective Displays.

\case{8.51} The Yugoslav Player is not eligible to receive guerrilla reinforcements (see 7.4) or Victory Points (see 14.0) when occupying any of the above Objective Displays unless the intrinsic Defense Strength of the Display has been eliminated (see 8.53).

\case{8.52} An intrinsic Defense Strength is equivalent to a normal Combat Strength, but it may never be used to attack. It may only be employed when the Yugoslav Player is conducting an attack against one of the above Displays. It may be employed by itself or in conjunction with Axis units currently occupying the Display. It may not be moved and may never be considered part of an Axis garrison.

\case{8.53} An intrinsic Defense Strength is considered eliminated in any attack against the Display that results in a ‘‘D’’ outcome. The elimination of a Display’s intrinsic Defense Strength must be kept track of on a separate piece of paper.

\case{8.54} The elimination of an intrinsic Defense Strength may never be used to satisfy Strength Point losses for Axis units that were affected by a ``D'' result in this Display.

\case{8.55} The intrinsic Defense Strengths of Trieste, Pola, Fiume, and Zara are not eliminated at the instant of Italian surrender (see 10.3).

\subsection{Combat Results Table}\label{sec:8.6}

\textbf{Addenda:} Modifications for Tito, Anti-Guerrilla Operations, and
mountain units are made to the probability ratio (8.21B) and are not column
shifts on the CRT. See the mapsheet.
